\documentclass[12pt]{amsbook}
\usepackage{preamble}

\begin{document}
\pagenumbering{gobble} % This kills the page numbering

\begin{center}
   \textsc{\large MATH 271, Exam 3}\\
\end{center}
\vspace{1cm}

\noindent\textbf{Name} \; \underline{\hspace{8cm}}

\vspace{1cm}

\noindent\textbf{Instructions} \; No textbook, homework, calculators, phones, or smart watches may be used for this exam. The exam is designed to take 50 minutes and must be submitted at the end of the class period. All of your solutions should be easily identifiable and supporting work must be shown. You may use any part of this packet as scratch paper, but please clearly label what work you want to be considered for grading. Ambiguous or illegible answers will not be counted as correct.\\

\noindent\emph{Only the highest scoring \underline{five} problems will be counted towards your total score. You cannot get over 75 points.}

\vspace{1cm}

\begin{flushleft}
\textbf{Problem 1} \; \underline{\hspace{1cm}}/15

\vspace{.25cm}

\textbf{Problem 2} \; \underline{\hspace{1cm}}/15

\vspace{.25cm}

\textbf{Problem 3} \; \underline{\hspace{1cm}}/15

\vspace{.25cm}

\textbf{Problem 4} \; \underline{\hspace{1cm}}/15

\vspace{.25cm}

\textbf{Problem 5} \; \underline{\hspace{1cm}}/15

\vspace{.25cm}

\textbf{Problem 6} \; \underline{\hspace{1cm}}/15

\vspace{.5cm}

\textbf{Total} \;\hspace{1.1cm} \underline{\hspace{1.25cm}}/75
\end{flushleft}

\vspace*{4cm}


\begin{center}\large{There are extra pages between each problem for scratch work.\\

Please circle your answers!}\end{center}









% Problem 1
\newpage
\begin{problem} \textbf{(3pts. Each)}\\

\def\arraystretch{4}% increase vertical spacing
\noindent\begin{tabularx}{\textwidth}{cXcc}
 & & T & F\\
(\theabc) & The cross product between two unit vectors is largest when the vectors are perpendicular. & \answerbox & \answerbox\\
(\theabc) & The inverse of the product matrix $[A][B]$ is $[B]^{-1}[A]^{-1}$. & \answerbox & \answerbox\\
(\theabc) & If we have the inhomogeneous equation
\[
[A]\vecv = \vecb,
\]
with $\vecb\neq \zerovec$, then $\det([A])=0$ means there are always infinitely many solutions. & \answerbox & \answerbox\\
(\theabc) & Every square $n\times n$-matrix is invertible. & \answerbox & \answerbox\\
(\theabc) & Hermitian matrices have real eigenvalues. & \answerbox & \answerbox\\
\end{tabularx}
\end{problem}

\newpage
\emph{Intentionally left blank to be used as scratch paper.}\\


% Problem 2
\newpage
\begin{problem}~
\begin{enumerate}[(a)]
    \item \textbf{(3pts.)} Draw and clearly label the vectors
\[
\vecv=\xhat + 3\yhat \qquad \vecw=4\xhat + 2\yhat
\]
in the plane provided using the provided grid.

\item \textbf{(4pts.)} Evaluate $\vecu=\vecv+\vecw$ and then draw and clearly label $\vecu$ in the plane.

\item \textbf{(4pts.)} Evaluate $\vecp=2\vecv-\sfrac{1}{2}\vecw$ and then draw and clearly label $\vecp$ in the plane.

\item \textbf{(4pts.)} Let $\vecr$ be the reflection of $\vecv$ across the $x$-axis. Draw and clearly label $\vecr$ in the plane.
\end{enumerate}

    \begin{center}
    \begin{tikzpicture}
    \draw[thin,gray!80] (-7,-7) grid (7,7);
    \draw[<->] (-7,0)--(7,0) node[right]{$x$};
    \draw[<->] (0,-7)--(0,7) node[above]{$y$};
    \end{tikzpicture}
    \end{center}

\end{problem}


\newpage
\emph{Intentionally left blank to be used as scratch paper.}\\


% Problem 3
\newpage
\begin{problem}
Consider the vectors 
\[
\vecu = \xhat + \yhat, \quad \vecv=2\zhat, \quad \vecw = \xhat -\yhat- \zhat.
\]
\begin{enumerate}[(a)]
    \item \textbf{(4pts.)} Why must the $z$-component of $\vecu \times \vecv$ be zero? Explain.
    \vspace*{5cm}
    \item \textbf{(6pts.)} Let $\vecr = 2\xhat +\alpha \yhat + 2\zhat$. For what value(s) of $\alpha$ is are the vectors $\vecw$ and $\vecr$ orthogonal?
    \vspace*{6cm}
    \item \textbf{(5pts.)} Compute the volume of the parallelepiped generated by the vectors $\vecu$, $\vecv$, and $\vecw$.
    \vspace*{6cm}
    
\end{enumerate}
\end{problem}

\newpage
\emph{Intentionally left blank to be used as scratch paper.}\\


% Problem 4
\newpage
\begin{problem} 
Let
\vspace*{.25cm}
\[
[M]= \begin{pmatrix} 1 & 2\\ 0 & 1 \end{pmatrix} \qquad 
[P]=\begin{pmatrix} 1 \\ 1 \end{pmatrix} \qquad 
[Q]=\begin{pmatrix} 2 & 1 \end{pmatrix} \qquad
[S]=\begin{pmatrix} 3 & 3 & 3 \\ 3 & 3 & 3 \end{pmatrix}
\]
\vspace*{.25cm}
\begin{enumerate}[(a)]
    \item \textbf{(5pts.)} Circle the matrix products can you compute.
    \vspace*{.75cm}
\[
[M][M], \qquad [P][P], \qquad [Q][P], \qquad [M][S], \qquad [S][M].
\]
\vspace*{.5cm}
    \item \textbf{(6pts.)} Compute the following:\\
    
    \begin{enumerate}[i.]
        \item $[A]=[P][Q]$;
        \item $[B]=[Q]^T[P]^T$. Is this equal to $([P][Q])^T$?
    \end{enumerate}
    \vspace*{7cm}
    \item \textbf{(4pts.)} Construct a $2\times 2$-matrix that swaps the $x$-component of a vector with the $y$-component of the vector.
\end{enumerate}

\end{problem}

\newpage
\emph{Intentionally left blank to be used as scratch paper.}\\


% Problem 5
\newpage
\begin{problem} 
Consider the system of linear equations:
\begin{align*}
    x + 2y &= 3\\
    x + y  &= 3
\end{align*}
\begin{enumerate}[(a)]
    \item \textbf{(3pts.)} Write this system in the form:
    \[
    [A]\vecv = \vecb.
    \]
    \vspace*{3cm}
    \item \textbf{(4pts.)} Argue that $[A]$ is invertible using the determinant.
    \vspace*{5cm}
    \item \textbf{(4pts.)} Compute the inverse of the matrix $[A]$.
    \vspace*{6cm}
    \item \textbf{(4pts.)} What is the solution $\vecv$ to the linear system?
\end{enumerate}
\end{problem}

\newpage
\emph{Intentionally left blank to be used as scratch paper.}\\


% Problem 6
\newpage
\begin{problem}
Consider the matrix
\[
[A]=\begin{pmatrix} 3 & -1\\ -1 & 3 \end{pmatrix}
\]
\begin{enumerate}[(a)]
    \item \textbf{(5pts.)} Show that $\evec_{1}=\xhat+\yhat$ and $\evec_2=-\xhat+\yhat$ are eigenvectors corresponding to the eigenvalues $\lambda_1=2$, and $\lambda_2=4$ respectively.
    \vspace*{4cm}
    \item \textbf{(3pts.)} Is the vector $\vecv=\evec_1+\evec_2$ also an eigenvector? Explain.
    \vspace*{4cm}
    \item \textbf{(3pts.)} Compute the trace and determinant of $[A]$ solely using the eigenvalues.
    \vspace{4cm}
    \item \textbf{(4pts.)} We can create diagonal matrix $[\Lambda]$ that is similar to $[A]$ by
    \[
    [\Lambda] = [P]^{-1}[A][P].
    \]
    Construct the matrix $[P]$ and $[\Lambda]$. \emph{Note, there is no need to construct $[P]^{-1}$ and multiply the three matrices above to find $[\Lambda]$!}
\end{enumerate}
\end{problem}

\newpage
\emph{Intentionally left blank to be used as scratch paper.}\\

\end{document}  